\chapter{Online Convex Optimization with Stochastic Energy Harvesting}
\label{Section:ChapEnergyharvest}

\noindent\textbf{Source and author contributions.}
This chapter is based on research published jointly with Michael J. Neely
in~\cite{10.1145/3428337} and presented at \emph{ACM SIGMETRICS 2021}.
Neely developed the main research idea and proposed most of the original
algorithm. I contributed algorithmic refinements, including separating
the allocation into total energy usage and allocation proportions.
I developed most of the proofs and wrote the initial manuscript.
Neely reviewed the proofs and manuscript and provided suggestions.

To simplify the presentation, this chapter specializes the general energy
bounds used in the published work~\cite{10.1145/3428337} to
nonnegative arrivals and total per-slot energy usage, both bounded above
by one.
The core usage-and-proportion algorithm is retained from the published
paper. Most of the subsequent changes in this chapter were produced with
AI assistance at my request. These include reorganizing the exposition
and proofs, making the conditional-mean assumption explicit, correcting
and completing the sufficient battery-capacity calculation, tightening
the initial-potential term in the regret bound, and making step-size
tuning explicit. The published paper's simulation study is not reproduced
here.

\section{Introduction}

Consider a system that allocates harvested energy among $n$ subsystems over
$T$ time slots, storing unused energy in a finite battery. On slot $t$, the
controller chooses a nonnegative vector
\[
    \mathbf{X}_t\coloneqq\left(X_{t,1},\ldots,X_{t,n}\right)^\top,
\]
where $X_{t,i}$ is the energy assigned to subsystem $i$. This decision has
a cost $c_t(\mathbf{X}_t)$. The function $c_t$ is convex, but it is unknown when $\mathbf{X}_t$ is
chosen and is revealed only after the decision.

The costs need not follow a probabilistic model. Energy arrivals are random,
and their distribution is unknown. Writing $B_t$ for the energy stored
at the start of slot $t$ and $A_t$ for the new arrival, the allocation must
satisfy
\begin{equation}\label{eq:battery-constraint}
    \sum_{i=1}^n X_{t,i}\leq B_t+A_t.
\end{equation}
Unused energy is stored up to the capacity $B_{\mathrm{max}}$, so
\begin{equation}\label{eq:B-update}
    B_{t+1}=\min\left\{B_t+A_t-\sum_{i=1}^n X_{t,i},
                    B_{\mathrm{max}}\right\},
    \qquad t\in[T].
\end{equation}
Here $B_1$ is the initial energy. Constraint \eqref{eq:battery-constraint}
ensures that the battery update remains nonnegative.

Choosing $\mathbf{X}_t$ before observing $c_t$ is the standard setting of
online convex optimization (OCO)
\cite{cesa1996worst,zinkevich2003online,hazan2016introduction,shalev2011online,bubeck2011introduction}.
The battery recursion couples allocations across slots.
Spending energy now leaves less for later slots, when it might be more useful.
Saving too much can leave energy unused or fill the battery, causing new
arrivals to be lost.

We measure performance using cumulative regret: the controller's total cost
minus the cost of the best nonnegative allocation $\mathbf{X}^*$ held fixed
across slots and chosen with knowledge of the full cost sequence. In addition
to the per-slot total-usage limit in Section~\ref{sec:energy:problem}, this
comparator must satisfy
\[
    \sum_{i=1}^nX_i^*\leq\overline{A},
    \qquad \overline{A}\coloneqq\expect{A_t}.
\]
This mean-budget comparator need not satisfy the battery constraint on
every slot.

The method in \cite{Yu2019} achieves an $O(\sqrt{T})$ regret bound with an
$O(\sqrt{T})$ battery capacity. Its regret bound contains a factor of
$O(\sqrt{n})$. This chapter reduces that factor to $O(\sqrt{\log n})$ for
$n\geq2$. For this comparison,
we keep the bound on each component of the cost subgradients fixed as
$n$ varies.

The available energy changes with arrivals and past spending, so a
simplex update alone cannot enforce the total energy budget. We separate
total usage from allocation proportions: the usage update enforces the
battery constraint, and entropic mirror descent updates the proportions
on a fixed probability simplex. The latter gives the dimension dependence
stated above.

\subsection{Wireless transmission example}\label{System example}

For $n$ parallel wireless channels, let $S_{t,i}\geq0$ be the power gain,
$N_{t,i}>0$ the noise power, and $W_i>0$ the effective bandwidth.
With unit-duration slots, additive white Gaussian noise, and channel
conditions fixed within each slot, the negative-throughput cost from
Shannon's formula~\cite[Section~5.2.1, Eq.~(5.12)]{tse2005fundamentals} is
\[
    c_t(\mathbf{X}_t)
    \coloneqq-\sum_{i=1}^n W_i\log_2\left(
       1+\frac{X_{t,i}S_{t,i}}{N_{t,i}}\right).
\]

This function is convex in the nonnegative allocation. The channel
parameters are unknown when the allocation is chosen and revealed
afterward; they need not be independent across channels or slots.
The energy arrivals must satisfy the assumptions in
Section~\ref{Assumptions}, including \eqref{eq:energy:conditional-mean}.
For solar harvesting, the i.i.d. arrival model is a simplification:
the solar trace studied in~\cite[Section~6.2]{10.1145/3428337}
exhibits approximately periodic behavior and is treated there as non-i.i.d.

For the uniform subgradient bound in \eqref{eq:G},
\[
 \left|\frac{\partial c_t}{\partial X_i}(\mathbf X)\right|
 =\frac{W_iS_{t,i}}{\ln 2\,(N_{t,i}+X_iS_{t,i})}
 \overset{\mathrm{(A)}}{\leq}\frac{W_iS_{t,i}}{N_{t,i}\ln 2},
\]
where (A) uses $X_i,S_{t,i}\geq0$ and $N_{t,i}>0$.
A uniform bound on $W_iS_{t,i}/N_{t,i}$ therefore suffices.

\subsection{Related work}

The fixed-set OCO benchmark and the convex regret bound are reviewed in
Section~\ref{sec:prelim:notation}. The common mirror
descent analysis is in Section~\ref{sec:intro:mirror-descent}.

Unlike the slotwise battery constraint and state recursion here,
work on long-term constraints considers fixed constraint functions
\cite{mahdavi2012trading,jenatton2016adaptive,titov2018mirror,JMLR:v21:16-494,Yuan2018}
or stochastic time-varying constraints
\cite{yu2017online,chen2018bandit,liakopoulos2019cautious}. Bregman methods
also appear in constrained OCO \cite{wei2020online}, although that model
does not include the physical storage dynamics studied here.

In more general online learning models with arbitrarily varying rewards
and sample-path average constraints, the best feasible performance in
hindsight need not be attainable \cite{mannor2009online}.
Here, the conditional-mean condition in
\eqref{eq:energy:conditional-mean} provides a predictable expected energy
supply given the past, while the costs need not follow a probabilistic
model.

The closest work, \cite{Yu2019}, combines online learning with finite
energy storage. Its regret and capacity bounds are compared above.
Packet-scheduling models minimize transmission completion time or maximize
finite-horizon throughput \cite{yang2011optimal,tutuncuoglu2012optimum}.
Other work learns policies with unknown Markov models for energy and data
arrivals \cite{blasco2013learning} or studies transmission under
time-correlated energy supply \cite{michelusi2013transmission}.
With i.i.d. arrivals, online power-control policies target long-term
throughput or general concave utility
\cite{shaviv2016universally,arafa2017energy}.
Online throughput maximization with battery overflow is studied
in \cite{wu2016online}; network models couple energy management with
traffic utility \cite{Huang2013,leonidas-energy-harvest}.
Here, convex costs are revealed only after allocation, and regret is
measured against a fixed allocation constrained by the mean energy arrival.

Finite-capacity inventory models also couple decisions through stored resources
\cite{yang2020online,lin2019competitive,chau2016cost,neely-inventory-control-cdc2010}.
For example, \cite{yang2020online} treats a finite-capacity inventory with a
linear objective that is known at the decision time. Here the objective
can be nonlinear and is revealed only after the decision.

\subsection{Contributions and scope}

We give explicit usage and proportion updates and a sufficient capacity
under which the available-energy cap does not change the usage decisions.
The tuned controller uses a planned horizon and the known
subgradient bound, but does not require the mean energy arrival.

\section{Problem formulation}
\label{sec:energy:problem}

We follow the conventions in Section~\ref{sec:intro:notation}. The system
has $n\geq2$ subsystems, and regret is evaluated over slots $t\in[T]$ for
an integer $T\geq1$. The recursions and their one-step and feasibility
bounds apply to any continuation with the same parameters and bounds.

The fixed allocation set is
\begin{equation}\label{eq:X-set}
\mathcal{X}\coloneqq\left\{\mathbf{x}\in\Real_+^n:
\sum_{i=1}^n x_i\leq1\right\}.
\end{equation}

Let $A_t$ denote the energy arriving on slot $t$. The battery has
capacity $B_{\mathrm{max}}$ and starts empty, $B_1=0$. At the beginning of
slot $t$, the controller observes the stored energy $B_t$ and the
new arrival $A_t$, then selects an allocation $\mathbf{X}_t$ satisfying
\begin{align}
\sum_{i=1}^n X_{t,i}&\leq B_t+A_t,
\label{eq:X-constraint1}\\
\mathbf{X}_t&\in\mathcal X.\label{eq:X-constraint2}
\end{align}
After the allocation is chosen, the cost function
$c_t:\mathcal X\to\Real$ is revealed. This function is continuous and convex.
The controller incurs cost $c_t(\mathbf{X}_t)$. The battery is then updated
according to \eqref{eq:B-update}.
The algorithm uses a subgradient of $c_t$ at $\mathbf{X}_t$.

\subsection{Assumptions}\label{Assumptions}

The arrivals are independent and follow the same unknown distribution.
Their mean is $\overline A\coloneqq\E[A_t]$. Each arrival satisfies
\begin{equation}\label{eq:energy:arrival-bound}
0\leq A_t\leq1\qquad\text{with probability one}.
\end{equation}
In particular, $0\leq\overline A\leq1$.

The costs may vary arbitrarily over time. The entire cost sequence may be
fixed in advance, or costs may be chosen during operation using observations
from earlier slots and random choices. For the expected regret guarantee,
we require that the history through slot $t$ does not change the expected
arrival on slot $t+1$.

Let $\mathcal F_t$ contain the history through the end of slot $t$. This
includes observed arrivals, allocations, costs, battery states, and any
random choices already used to generate them. The requirement is
\begin{equation}\label{eq:energy:conditional-mean}
\E[A_{t+1}\mid\mathcal F_t]=\overline A.
\end{equation}
For example, this condition holds when the arrivals are independent with
the same distribution, are independent of the random choices used to select
costs, and costs are chosen without access to future arrivals.

The condition concerns the whole history because earlier costs could
reveal information about future arrivals. Requiring only $A_t$ and $c_t$
to be independent on the same slot would not rule this out.

At every slot $t$ and every point $\mathbf{x}\in\mathcal X$, assume that the
controller can obtain a subgradient
$\mathbf{g}_t(\mathbf{x})\in\partial c_t(\mathbf{x})$ satisfying
\begin{equation}\label{eq:G}
\|\mathbf{g}_t(\mathbf{x})\|_\infty
=\max_{1\leq i\leq n}|g_{t,i}(\mathbf{x})|\leq G.
\end{equation}
Here, $G>0$ is a known constant.
In particular, the subgradient inequality gives
\begin{equation}\label{eq:energy:subgradient}
c_t(\mathbf{x})-c_t(\mathbf{y})\leq \mathbf{g}_t(\mathbf{x})^\top(\mathbf{x}-\mathbf{y})
\qquad(\mathbf{x},\mathbf{y}\in\mathcal X).
\end{equation}
Lemma~\ref{lem:intro:lipschitz}, with $Q=\mathcal X$, the $\ell_1$ norm,
its dual $\ell_\infty$ norm, and $L=G$ from \eqref{eq:G}, also gives
\begin{equation}\label{eq:energy:lipschitz}
|c_t(\mathbf{x})-c_t(\mathbf{y})|\leq G\|\mathbf{x}-\mathbf{y}\|_1.
\end{equation}
The bound in \eqref{eq:G} applies to the subgradient the controller selects
at each point. It need not hold for every possible subgradient.

\subsection{Comparison allocation and regret}

We compare the algorithm with an allocation that stays the same on every
slot. Its total usage cannot exceed the mean arrival
$\overline{A}=\expect{A_t}$. This fixed allocation need not satisfy the
battery constraint on every slot. It must satisfy
\begin{align}
  U^*\coloneqq\sum_{i=1}^n X_i^*&\leq\overline{A}, \label{eq:A1}\\
  \mathbf{X}^*&\in\mathcal{X}. \label{eq:A2}
\end{align}
Since $\overline A\leq1$, the allocations satisfying
\eqref{eq:A1}--\eqref{eq:A2} form the set
\begin{equation}\label{eq:energy:comparison-set}
  \mathcal A\coloneqq\left\{\mathbf x\in\Real_+^n:
    \sum_{i=1}^n x_i\leq\overline A\right\}.
\end{equation}
The set $\mathcal A$ contains the zero allocation and is compact,
and the costs are continuous. A best fixed allocation therefore exists
for every realized cost sequence. We define regret by
\begin{equation}
\label{eq:ch3-regret-definition}
  \Reg_T(\mathcal A)
  \coloneqq\sum_{t=1}^T c_t(\mathbf{X}_t)
    -\min_{\mathbf{X}\in\mathcal{A}}\sum_{t=1}^T c_t(\mathbf{X}).
\end{equation}
The minimizing vector may depend on the full cost sequence, including when
costs depend on past observations.

\subsection{Main bounds}\label{sec:energy:main-bounds}

Under the assumptions of Section~\ref{Assumptions}, fix a planned horizon
$T\geq1$. Algorithm~\ref{Alg1}, initialized as specified below and using
\eqref{eq:ch3-parameter-scaling} and \eqref{eq:lambda-opt}, satisfies
\[
 B_{\mathrm{max}}=3\sqrt T,\qquad
 \E[\Reg_T(\mathcal A)]
 \leq G\sqrt T\left[6+\sqrt{2\log n}\right].
\]
The capacity choice makes the relaxed usage update feasible on every slot
by Theorem~\ref{Theorem of battery}. The regret bound follows from
Theorem~\ref{Main Theorem} and \eqref{eq:ch3-explicit-sqrt-bound}.
For fixed $G$, it is $O(\sqrt{T(1+\log n)})$, and the sufficient capacity
is independent of $n$.

\section{Algorithm}
\label{ALGorithm}

\subsection{Energy usage and allocation proportions}

Write an allocation as
\begin{equation}\label{eq:energy:decomposition}
\mathbf{X}_t=U_t\mathbf{P}_t,\qquad
U_t\coloneqq\sum_{i=1}^n X_{t,i},\qquad \mathbf{P}_t\in\Delta_n,
\end{equation}
where $\Delta_n$ is the probability simplex in~\eqref{eq:intro:simplices}.
The scalar $U_t$ is total usage, and $\mathbf P_t$ gives the allocation
proportions.
We continue updating the proportions even on slots when no energy is used.
The constraints become
\[
\mathbf P_t\in\Delta_n,\qquad
U_t\in[0,\min\{1,B_t+A_t\}].
\]
Write the observed cost subgradient and the corresponding usage subgradient as
\begin{equation}\label{eq:energy:gh}
\mathbf g_t\coloneqq\mathbf g_t(\mathbf X_t),\qquad
h_t\coloneqq\mathbf g_t^\top\mathbf P_t.
\end{equation}
By Lemma~\ref{lem:intro:convex-comparison}, $h_t$ is a subgradient with
respect to total usage when the proportions are fixed.
Since $\mathbf P_t$ is a probability vector,
\begin{equation}\label{eq:energy:h-bound}
|h_t|\leq\sum_{i=1}^nP_{t,i}|g_{t,i}|
\overset{\mathrm{(A)}}{\leq}G.
\end{equation}
Here (A) uses \eqref{eq:G} and $\sum_iP_{t,i}=1$; the first inequality
is the triangle inequality applied to \eqref{eq:energy:gh}.

Although $c_t$ is convex in $\mathbf X$, the function $c_t(u\mathbf p)$
need not be jointly convex in $(u,\mathbf p)$. It is convex in either
variable with the other fixed, by affine composition
\cite[Sections~3.1.4 and~3.2.2]{boyd_vandenberghe_2004}.
For a fixed comparator $\mathbf X^*=U^*\mathbf P^*\in\mathcal A$, take
$U^*=\sum_iX_i^*$ and $\mathbf P^*\in\Delta_n$; when $U^*=0$,
any $\mathbf P^*\in\Delta_n$ gives this factorization. Then
\begin{equation}\label{eq:energy:split-identity}
\mathbf X_t-\mathbf X^*
=(U_t-U^*)\mathbf P_t+U^*(\mathbf P_t-\mathbf P^*).
\end{equation}
Applying \eqref{eq:energy:subgradient} and substituting
\eqref{eq:energy:split-identity} gives
\begin{equation}\label{eq:energy:split-bound}
 c_t(\mathbf X_t)-c_t(\mathbf X^*) \overset{\mathrm{(A)}}{\leq}\mathbf g_t^\top(\mathbf X_t-\mathbf
   X^*) =h_t(U_t-U^*)+U^*\mathbf g_t^\top(\mathbf P_t-\mathbf P^*).
\end{equation}
Here (A) is the subgradient inequality at $\mathbf X_t$; the equality
uses \eqref{eq:energy:gh}. The usage and proportion updates use $h_t$
and $\mathbf g_t$, respectively. The analysis multiplies the proportion
bound by $U^*$, which is fixed across slots; neither update uses $U^*$.
Direct linearization in $\mathbf p$ at fixed $U_t$ would instead give
$U_t\mathbf g_t$.

\subsection{Two Bregman divergences}

We use separate Bregman proximity terms for changes in total usage and allocation
proportions. For usage, use the quadratic term $(u-v)^2/2$ from the
scalar case of \eqref{eq:intro:quadratic-geometry}.
For proportions, use the Kullback--Leibler (KL) divergence, obtained by
restricting the generalized KL divergence in~\eqref{eq:intro:generalized-kl}
to the probability simplex:
\begin{equation}\label{eq:energy:entropy-divergence}
D_\Psi(\mathbf p,\mathbf q)=\sum_{i=1}^n p_i\log\frac{p_i}{q_i},
\qquad \mathbf p,\mathbf q\in\Delta_n,\quad q_i>0.
\end{equation}
The linear terms in \eqref{eq:intro:generalized-kl} cancel because
$\sum_i p_i=\sum_i q_i=1$.
We use the convention $0\log0=0$. The initial proportions are uniform.
Lemma~\ref{lem:intro:entropy} therefore gives
\[
D_\Psi(\mathbf p,\mathbf P_1)\leq\log n
\qquad\text{for every }\mathbf p\in\Delta_n.
\]
These Bregman proximity terms give the mirror-descent updates
\cite{nemirovskij1983problem,BECK2003167,hazan2010extracting,cesa2006prediction}.

\subsection{Battery feedback and update rules}

Choose positive step sizes $\eta$ and $\lambda$ for usage and proportions,
and a price scale $\gamma>0$. At the end of slot $t$, the controller knows
the battery level $B_{t+1}$ and the cost subgradient $\mathbf g_t$. It uses
$\mathbf g_t$ to compute $\mathbf P_{t+1}$ and $B_{t+1}$ to compute an energy
price $\pi_t$. At the start of the next slot, it observes $A_{t+1}$ and then
chooses $U_{t+1}$ using this price and the available energy.

Define the battery price by
\begin{equation}\label{eq:energy:price}
\pi_t\coloneqq\gamma(B_{\mathrm{max}}-B_{t+1})\geq0.
\end{equation}
The price measures the deficit from full capacity: a lower battery level
gives a higher price and discourages usage, while a full battery has zero
price. With the usage step size fixed, $\gamma$ controls how strongly this
deficit affects the usage decision.
The battery recursion gives
\begin{equation}\label{eq:energy:price-update}
\pi_{t+1}=[\pi_t+\gamma(U_{t+1}-A_{t+1})]_+,
\qquad \pi_0=\gamma B_{\mathrm{max}}.
\end{equation}
Indeed, \eqref{eq:B-update} and
$\sum_iX_{t+1,i}=U_{t+1}$ from \eqref{eq:energy:decomposition} give
$B_{t+2}=\min\{B_{t+1}+A_{t+1}-U_{t+1},B_{\mathrm{max}}\}$.
Substitute this expression into \eqref{eq:energy:price}, and use
$b-\min\{z,b\}=\max\{b-z,0\}$ and $\gamma>0$.
Thus the price increases when usage exceeds arrivals and decreases, down
to zero, when arrivals exceed usage.

This battery-dependent penalty is related to drift-plus-penalty
energy-harvesting control in
\cite[Section~III-B and Lemma~3]{Yu2019}.
For the battery price $\pi_t$, define the quadratic potential
\begin{equation}\label{eq:energy:price-potential}
V_t\coloneqq\frac{\pi_t^2}{2\gamma}
   =\frac{\gamma}{2}(B_{\mathrm{max}}-B_{t+1})^2.
\end{equation}
Its one-slot change, or drift, satisfies
\[
V_{t+1}-V_t
\leq\pi_t(U_{t+1}-A_{t+1})
    +\frac{\gamma}{2}(U_{t+1}-A_{t+1})^2.
\]
Section~\ref{sec:energy:battery-potential} derives this inequality and
bounds the squared term by one. With $\pi_t$ fixed and $A_{t+1}$ observed,
the remaining usage-dependent term is $\pi_tU_{t+1}$.

The usage update is
\begin{equation}\label{eq:a-decision}
U_{t+1}\coloneqq\argmin_{u\in I_{t+1}}
\left\{\eta(h_t+\pi_t)u+\frac12(u-U_t)^2\right\},
\end{equation}
where
\[
I_{t+1}\coloneqq[0,\min\{1,B_{t+1}+A_{t+1}\}].
\]
This interval is nonempty because $B_{t+1}$ and $A_{t+1}$ are nonnegative.
Completing the square in \eqref{eq:a-decision} gives
$\tfrac12[u-(U_t-\eta(h_t+\pi_t))]^2$ plus a term independent of $u$.
Thus the usage update is the Euclidean projection onto $I_{t+1}$:
\begin{equation}\label{Update A}
U_{t+1}=[U_t-\eta(h_t+\pi_t)]
 _{0}^{\min\{1,B_{t+1}+A_{t+1}\}},
\end{equation}
where $[z]_\ell^r\coloneqq\min\{r,\max\{\ell,z\}\}$ for $\ell\leq r$.
Before projection, usage increases when the linearized marginal benefit
$-h_t$ exceeds the battery price $\pi_t$, and decreases when it falls below
the price. The projection enforces the usage limits and available energy.
The cap uses $B_{t+1}+A_{t+1}$, while $\pi_t$ uses $B_{t+1}$.

The proportion update is
\begin{equation}\label{eq:p-decision}
\mathbf P_{t+1}\coloneqq\argmin_{\mathbf p\in\Delta_n}
\{\lambda\mathbf g_t^\top\mathbf p+D_\Psi(\mathbf p,\mathbf P_t)\}.
\end{equation}
Lemma~\ref{lem:intro:entropy-update}, with $Q=\Delta_n$,
$q=\mathbf P_t$, $g=\mathbf g_t$, and its step size set to $\lambda>0$, yields
\begin{equation}\label{eq:p-decision-observe}
P_{t+1,i}
=\frac{P_{t,i}\exp(-\lambda g_{t,i})}
       {\sum_jP_{t,j}\exp(-\lambda g_{t,j})}.
\end{equation}
The objective in~\eqref{eq:p-decision} is strictly convex, so the proportion
update has a unique solution. The initial proportions are positive.
The update formula therefore keeps $P_{t,i}>0$ for every slot and subsystem.
Subsystems with smaller subgradient components gain proportion relative
to those with larger components.

\subsection{Algorithm specification}\label{section:alg-spec}

\begin{algorithm}
\caption{Energy usage, allocation proportions, and battery price}
\label{Alg1}
\begin{algorithmic}[1]
\Require Arrivals satisfying \eqref{eq:energy:arrival-bound}
\Require Battery capacity $B_{\mathrm{max}}>0$
\Require Step sizes $\eta,\lambda>0$ and price scale $\gamma>0$
\State $B_1\gets0$, $\pi_0\gets\gamma B_{\mathrm{max}}$,
       $U_1\gets 0$
\State $\mathbf P_1\gets(1/n,\ldots,1/n)^\top$
\For{$t=1,\ldots,T$}
  \State Observe $A_t$ and the stored energy $B_t$
  \If{$t\geq2$}
    \State $U_t\gets[U_{t-1}-\eta(h_{t-1}+\pi_{t-1})]
       _{0}^{\min\{1,B_t+A_t\}}$
  \EndIf
  \State Allocate $\mathbf X_t\gets U_t\mathbf P_t$
  \State Observe $c_t$, incur $c_t(\mathbf X_t)$, and obtain
         $\mathbf g_t\gets\mathbf g_t(\mathbf X_t)$
  \State $h_t\gets\mathbf g_t^\top\mathbf P_t$
  \State $B_{t+1}\gets\min\{B_t+A_t-U_t,B_{\mathrm{max}}\}$
  \State $\pi_t\gets\gamma(B_{\mathrm{max}}-B_{t+1})$
  \State Update $\mathbf P_{t+1}$ using \eqref{eq:p-decision-observe}
\EndFor
\end{algorithmic}
\end{algorithm}

Each slot uses $O(n)$ arithmetic operations, apart from obtaining the
subgradient. The algorithm needs $O(n)$ memory.

\section{Feasibility of the relaxed usage update}
\label{sec:energy:feasibility}

Using the fixed interval $[0,1]$ allows comparison with a fixed usage that
may exceed the available energy on individual slots. We show that
sufficient battery capacity makes the relaxed algorithm feasible, so its
regret bound also applies to Algorithm~\ref{Alg1}.

Fix $\eta,\gamma,\lambda>0$ and any sequence satisfying
\eqref{eq:energy:arrival-bound} and \eqref{eq:G}.
The relaxed algorithm replaces \eqref{eq:a-decision} by
\begin{equation}
\label{eq:a-decision-relax}
    U_{t+1}\coloneqq\argmin_{u\in[0,1]}
    \left\{\eta(h_t+\pi_t)u+\frac12(u-U_t)^2\right\},
\end{equation}
and retains the other updates and initializations of Algorithm~\ref{Alg1}.
Completing the square in \eqref{eq:a-decision-relax} gives its unique minimizer:
\begin{equation}
\label{eq:relaxed-A}
    U_{t+1}=
    \left[U_t-\eta(h_t+\pi_t)\right]
    ^{1}_{0}.
\end{equation}
\subsection{A lower bound on the battery reserve}

We prove the invariant $B_{t+1}\geq rU_t$. The analysis parameter $r\geq0$
will be chosen to minimize the sufficient capacity bound
\begin{equation}
\label{eq:Bmax-a}
    B_{\mathrm{req}}(r)
    \coloneqq\frac G\gamma+1+r+\frac{1}{\eta\gamma(1+r)}.
\end{equation}
\begin{Thm:Theorem}[Sufficient capacity and feasibility]
\label{Theorem of battery}
If $B_{\mathrm{max}}\geq B_{\mathrm{req}}(r)$ in \eqref{eq:Bmax-a},
$B_1=0$, and $U_1=0$, then the relaxed algorithm satisfies
\begin{equation}
\label{eq:induction}
    0\leq rU_t\leq B_{t+1}\leq B_{\mathrm{max}}
    \qquad\text{for every }t\geq1.
\end{equation}
Consequently,
\begin{equation}\label{eq:energy:relaxed-feasibility}
 U_t\leq B_t+A_t,\qquad t\geq1,
\end{equation}
and the relaxed algorithm coincides with Algorithm~\ref{Alg1}.
\end{Thm:Theorem}
\begin{proof}
By $U_1=0$ and \eqref{eq:relaxed-A}, $0\leq U_t\leq1$;
hence $rU_t\geq0$.
Also, \eqref{eq:B-update} gives $B_{t+1}\leq B_{\mathrm{max}}$.
For $t=1$,
\begin{equation}\label{eq:energy:reserve-base}
\begin{aligned}
 B_2&=\min\{A_1,B_{\mathrm{max}}\}\\
    &\geq0\\
    &=rU_1.
\end{aligned}
\end{equation}
This uses $B_1=U_1=0$, \eqref{eq:B-update}, and $A_1,B_{\mathrm{max}}\geq0$.
Assume \eqref{eq:induction} at slot $t$.

\emph{Case 1: $B_{t+1}\geq1+r$.} Then
\begin{equation}\label{eq:energy:reserve-high}
\begin{aligned}
 (1+r)U_{t+1}
 &\overset{\mathrm{(A)}}{\leq}1+r\\
 &\overset{\mathrm{(B)}}{\leq}B_{t+1}\\
 &\overset{\mathrm{(C)}}{\leq}B_{t+1}+A_{t+1}.
\end{aligned}
\end{equation}
Steps (A)--(C) use $U_{t+1}\leq1$, the case assumption, and $A_{t+1}\geq0$.

\emph{Case 2: $B_{t+1}<1+r$.} By \eqref{eq:energy:price},
\begin{equation}\label{eq:energy:low-battery-estimates}
\begin{aligned}
 h_t+\pi_t
 &\overset{\mathrm{(A)}}{\geq}-G+\pi_t\\
 &=-G+\gamma(B_{\mathrm{max}}-B_{t+1})\\
 &\overset{\mathrm{(B)}}{\geq}-G+\gamma(B_{\mathrm{req}}(r)-B_{t+1})\\
 &=-G+\gamma\left(\frac G\gamma+1+r
              +\frac{1}{\eta\gamma(1+r)}-B_{t+1}\right)\\
 &=-G+G+\gamma(1+r-B_{t+1})+\frac{1}{\eta(1+r)}\\
 &=\gamma(1+r-B_{t+1})+\frac{1}{\eta(1+r)}\\
 &\overset{\mathrm{(C)}}{\geq}\frac{1}{\eta(1+r)}.
\end{aligned}
\end{equation}
Steps (A)--(C) use \eqref{eq:energy:h-bound},
$B_{\mathrm{max}}\geq B_{\mathrm{req}}(r)$ with $\gamma>0$, and
$B_{t+1}<1+r$, respectively; the equalities substitute \eqref{eq:Bmax-a}.
Thus
\begin{equation}\label{eq:battery-usage-decrease}
\begin{aligned}
 U_{t+1}
 &=[U_t-\eta(h_t+\pi_t)]_0^1\\
 &\overset{\mathrm{(A)}}{\leq}
       \left[U_t-\frac{\eta}{\eta(1+r)}\right]_0^1\\
 &=\left[U_t-\frac{1}{1+r}\right]_0^1\\
 &=\max\left\{U_t-\frac{1}{1+r},0\right\}.
\end{aligned}
\end{equation}
Here (A) uses the preceding bound and monotonicity of clipping with $\eta>0$.
The equalities use \eqref{eq:relaxed-A}, $U_t\leq1$, and $1/(1+r)>0$. Thus
\begin{equation}\label{eq:energy:reserve-low}
\begin{aligned}
 (1+r)U_{t+1}
 &\overset{\mathrm{(A)}}{\leq}
      (1+r)\max\left\{U_t-\frac{1}{1+r},0\right\}\\
 &=\max\left\{(1+r)\left(U_t-\frac{1}{1+r}\right),0\right\}\\
 &=\max\{(1+r)U_t-1,0\}\\
 &=\max\{rU_t+U_t-1,0\}\\
 &\overset{\mathrm{(B)}}{\leq}\max\{rU_t,0\}\\
 &=rU_t\\
 &\overset{\mathrm{(C)}}{\leq}B_{t+1}\\
 &\overset{\mathrm{(D)}}{\leq}B_{t+1}+A_{t+1}.
\end{aligned}
\end{equation}
Step (A) uses \eqref{eq:battery-usage-decrease} and $1+r>0$;
(B) uses $U_t\leq1$; (C) is the induction hypothesis; and (D) uses
$A_{t+1}\geq0$. The maximum is monotone, and $rU_t\geq0$.

In both cases, \eqref{eq:B-update} gives
\begin{equation}\label{eq:energy:reserve-step}
\begin{aligned}
 B_{t+2}-rU_{t+1}
 &=\min\{B_{t+1}+A_{t+1}-U_{t+1},B_{\mathrm{max}}\}-rU_{t+1}\\
 &=\min\{B_{t+1}+A_{t+1}-U_{t+1}-rU_{t+1},B_{\mathrm{max}}-rU_{t+1}\}\\
 &=\min\{B_{t+1}+A_{t+1}-(1+r)U_{t+1},B_{\mathrm{max}}-rU_{t+1}\}\\
 &\overset{\mathrm{(A)}}{\geq}\min\{0,B_{\mathrm{max}}-rU_{t+1}\}\\
 &\overset{\mathrm{(B)}}{\geq}\min\{0,B_{\mathrm{max}}-r\}\\
 &\overset{\mathrm{(C)}}{\geq}\min\{0,0\}\\
 &=0.
\end{aligned}
\end{equation}
By monotonicity of the minimum, (A) follows from
\eqref{eq:energy:reserve-high} or \eqref{eq:energy:reserve-low};
(B) uses $r\geq0$ and $U_{t+1}\leq1$; and (C) uses
$B_{\mathrm{max}}\geq B_{\mathrm{req}}(r)>r$ from \eqref{eq:Bmax-a}.
This completes the induction. Finally,
\begin{equation}\label{eq:energy:feasibility-consequence}
\begin{aligned}
 U_t&\overset{\mathrm{(A)}}{\leq}U_t+B_{t+1}\\
 &\overset{\mathrm{(B)}}{\leq}U_t+(B_t+A_t-U_t)\\
 &=B_t+A_t.
\end{aligned}
\end{equation}
Here (A) uses \eqref{eq:induction}; (B) uses \eqref{eq:B-update}.
Thus \eqref{eq:energy:relaxed-feasibility} holds, and the minimizer in
\eqref{eq:a-decision-relax} belongs to the interval in \eqref{eq:a-decision}.
Uniqueness and identical initialization give identical algorithms.
\end{proof}

\subsection{Minimizing the sufficient battery-capacity bound}

The following calculation corrects the capacity tuning in
\cite[Section~4.2]{10.1145/3428337}, under the unit bounds used here.

For $0<\eta\gamma\leq1$, \eqref{eq:Bmax-a} gives
\begin{equation}\label{eq:energy:capacity-square}
\begin{aligned}
 B_{\mathrm{req}}(r)-\frac G\gamma-\frac{2}{\sqrt{\eta\gamma}}
 &=1+r+\frac{1}{\eta\gamma(1+r)}-\frac{2}{\sqrt{\eta\gamma}}\\
 &=\left(\sqrt{1+r}-\frac{1}{\sqrt{\eta\gamma(1+r)}}\right)^2\\
 &\geq0.
\end{aligned}
\end{equation}
The inequality in \eqref{eq:energy:capacity-square} is nonnegativity
of a square. Equality holds at $r=1/\sqrt{\eta\gamma}-1\geq0$,
where nonnegativity follows from $\eta\gamma\leq1$.
For $\eta\gamma>1$,
\begin{equation}\label{eq:energy:capacity-boundary}
\begin{aligned}
 B_{\mathrm{req}}(r)-B_{\mathrm{req}}(0)
 &=r+\frac{1}{\eta\gamma(1+r)}-\frac{1}{\eta\gamma}\\
 &=r-\frac{r}{\eta\gamma(1+r)}\\
 &=r\left(1-\frac{1}{\eta\gamma(1+r)}\right)\\
 &\overset{\mathrm{(A)}}{\geq}r\left(1-\frac{1}{\eta\gamma}\right)\\
 &\overset{\mathrm{(B)}}{\geq}0.
\end{aligned}
\end{equation}
In \eqref{eq:energy:capacity-boundary}, (A) uses $1+r\geq1$ and $r\geq0$;
(B) uses $\eta\gamma>1$ and $r\geq0$. Equality holds at $r=0$.
Hence \eqref{eq:energy:capacity-square}--\eqref{eq:energy:capacity-boundary} give
\begin{equation}
\label{eq:battery-optimal-a}
    r^*\coloneqq\left[\frac{1}{\sqrt{\eta\gamma}}-1\right]_+.
\end{equation}
Substituting \eqref{eq:battery-optimal-a} into \eqref{eq:Bmax-a} yields
\begin{equation}
\label{B^*max}
\begin{aligned}
    B_{\mathrm{req}}^\star&\coloneqq\min_{r\geq0}B_{\mathrm{req}}(r)\\
    &=\begin{cases}
       \displaystyle\frac G\gamma+\frac{2}{\sqrt{\eta\gamma}},
       &0<\eta\gamma\leq1,\\[1.2ex]
       \displaystyle\frac G\gamma+1+\frac{1}{\eta\gamma},
       &\eta\gamma>1.
     \end{cases}
\end{aligned}
\end{equation}
Thus $B_{\mathrm{max}}\geq B_{\mathrm{req}}^\star$ in \eqref{B^*max}
suffices for Theorem~\ref{Theorem of battery} with $r=r^*$.

\section{Regret analysis of the relaxed algorithm}
\label{regert sec}

We use the comparison set $\mathcal A$ of allocations
$\mathbf X\in\mathcal X$ with $\sum_iX_i\leq\overline A$,
specified by \eqref{eq:A1}--\eqref{eq:A2}, and the regret defined in
\eqref{eq:ch3-regret-definition}.

Throughout this section, assume that $B_{\mathrm{max}}\geq B_{\mathrm{req}}(r)$
in \eqref{eq:Bmax-a} for some $r\geq0$. Under this condition,
Theorem~\ref{Theorem of battery} ensures that Algorithm~\ref{Alg1} uses
the relaxed update \eqref{eq:a-decision-relax}.

The price and subgradients in
\eqref{eq:energy:price} and \eqref{eq:energy:gh} satisfy
\begin{equation}
\label{eq:ch3-shifted-battery}
  0\leq\pi_t\leq\gamma B_{\mathrm{max}},\qquad
  \|\mathbf{g}_t\|_\infty\leq G,\qquad |h_t|\leq G.
\end{equation}
The price bound follows from \eqref{eq:energy:price} and
\eqref{eq:induction}; the other two bounds are \eqref{eq:G} and
\eqref{eq:energy:h-bound}.
For the proof only, extend the usage and battery recursions to slot $T+1$
using the next arrival $A_{T+1}$ under the same arrival model. The resulting
values $U_{T+1}$ and $B_{T+2}$ let us sum the one-step bounds through slot
$T$. No decision through slot $T$ uses this future arrival, and the regret
in \eqref{eq:ch3-regret-definition} includes costs only through slot $T$.

Adding the usage bound, $U^*$ times the proportion bound, and the price
drift cancels $\pi_tU_{t+1}$. The remaining term is bounded by
$\pi_t(\overline A-A_{t+1})$, whose conditional mean is zero by
\eqref{eq:energy:conditional-mean}.

\subsection{Bounds for the two mirror-descent updates}

\begin{Thm:Lemma}[Updating allocation proportions]
\label{lem:ch3-direction-bound}
For every $\mathbf{P}^*\in\Delta_n$ and every $t\geq1$,
\begin{equation}
\label{eq PArt1}
  \mathbf{g}_t^\top(\mathbf{P}_t-\mathbf{P}^*)
  \leq\frac{\lambda G^2}{2}
    +\frac{D_\Psi(\mathbf{P}^*,\mathbf{P}_t)-D_\Psi(\mathbf{P}^*,\mathbf{P}_{t+1})}{\lambda}.
\end{equation}
\end{Thm:Lemma}

\begin{proof}
Apply Lemma~\ref{lem:intro:mirror-descent} to \eqref{eq:p-decision}
with $Q=\Delta_n$, $\Phi=\Psi$, $\sigma=1$, step size $\lambda>0$,
and $(x,x^+,u)=(\mathbf P_t,\mathbf P_{t+1},\mathbf P^*)$ in the
$\ell_1/\ell_\infty$ norm pair. The curvature follows from
\eqref{eq:intro:entropy-curvature}; the reference points are positive by
\eqref{eq:p-decision-observe}, while the comparator may lie on the boundary:
\[
\begin{aligned}
 \mathbf g_t^\top(\mathbf P_t-\mathbf P^*)
 &\overset{\mathrm{(A)}}{\leq}
 \frac{D_\Psi(\mathbf P^*,\mathbf P_t)-D_\Psi(\mathbf P^*,\mathbf P_{t+1})}{\lambda}
 +\frac{\lambda}{2}\|\mathbf g_t\|_\infty^2\\
 &\overset{\mathrm{(B)}}{\leq}
 \frac{D_\Psi(\mathbf P^*,\mathbf P_t)-D_\Psi(\mathbf P^*,\mathbf P_{t+1})}{\lambda}
 +\frac{\lambda G^2}{2}.
\end{aligned}
\]
Step (A) is \eqref{eq:intro:md-one-step}; (B) uses \eqref{eq:G}.
\end{proof}

\begin{Thm:Lemma}[Usage update]
\label{lem:ch3-usage-bound}
For every $U^*\in[0,1]$ and every $t\geq1$,
\begin{equation}
\label{eq Part2}
 h_t(U_t-U^*)+\pi_t(U_{t+1}-U^*) \leq{}\frac{\eta G^2}{2}
   +\frac{(U^*-U_t)^2-(U^*-U_{t+1})^2}{2\eta}.
\end{equation}
\end{Thm:Lemma}

\begin{proof}
Apply \eqref{eq:pushback} in Lemma~\ref{lem:pushback} to
\eqref{eq:a-decision-relax} with $Q=[0,1]$,
$f(u)=\eta(h_t+\pi_t)u$, $\alpha=1$, and
$(y,\hat x,z)=(U_t,U_{t+1},U^*)$, using the quadratic generator $u^2/2$.
Both reference points belong to $Q$, and the generator is differentiable
on $\Real$. Thus
\begin{equation}
\label{eq:main2}
\begin{aligned}
 &\eta h_t(U_t-U^*)+\eta\pi_t(U_{t+1}-U^*)\\
 &\quad=\eta(h_t+\pi_t)(U_{t+1}-U^*)+\eta h_t(U_t-U_{t+1})\\
 &\quad\overset{\mathrm{(A)}}{\leq}
   \frac{(U^*-U_t)^2-(U^*-U_{t+1})^2}{2}
   +\eta h_t(U_t-U_{t+1})-\frac12(U_t-U_{t+1})^2\\
 &\quad=\frac{(U^*-U_t)^2-(U^*-U_{t+1})^2}{2}
   +\frac{\eta^2h_t^2}{2}-\frac12(U_t-U_{t+1}-\eta h_t)^2\\
 &\quad\overset{\mathrm{(B)}}{\leq}
   \frac{(U^*-U_t)^2-(U^*-U_{t+1})^2}{2}+\frac{\eta^2G^2}{2}.
\end{aligned}
\end{equation}
Step (A) substitutes the quadratic divergence into \eqref{eq:pushback}.
The next equality completes the square. Step (B) drops that square and
uses $|h_t|\leq G$ from \eqref{eq:energy:h-bound}.
Divide by $\eta>0$ to obtain \eqref{eq Part2}.
\end{proof}
The price term uses $U_{t+1}$. This keeps the price out of the error term,
which is $\eta G^2/2$.

\subsection{A common learning and battery potential}
\label{sec:energy:battery-potential}

For $V_t=\pi_t^2/(2\gamma)$ in \eqref{eq:energy:price-potential},
the price recursion \eqref{eq:energy:price-update} and
$[z]_+^2\leq z^2$ give
\begin{equation}\label{eq:include}
\begin{aligned}
 V_{t+1}-V_t
 &\overset{\mathrm{(A)}}{\leq}\pi_t(U_{t+1}-A_{t+1})
   +\frac\gamma2(U_{t+1}-A_{t+1})^2\\
 &\overset{\mathrm{(B)}}{\leq}\pi_t(U_{t+1}-A_{t+1})+\frac{\gamma}{2}.
\end{aligned}
\end{equation}
Step (A) expands $[\pi_t+\gamma(U_{t+1}-A_{t+1})]^2-\pi_t^2$
and divides by $2\gamma>0$. For (B), both $A_{t+1}$ and $U_{t+1}$
belong to $[0,1]$, so $(A_{t+1}-U_{t+1})^2\leq1$.

\begin{Thm:Lemma}[Comparison using the mean energy arrival]
\label{lem:ch3-mean-energy}
For every comparison usage $U^*\leq\overline{A}$, we have
\[
  \pi_t(U^*-A_{t+1})\leq\pi_t(\overline A-A_{t+1})
\]
for every realization. This also holds when $U^*$ is random and chosen
after observing the full cost sequence. The expression on the right has
conditional mean zero:
\begin{equation}
\label{quu1 bound}
  \expect{\pi_t(\overline A-A_{t+1})\mid\mathcal F_t}=0.
\end{equation}
\end{Thm:Lemma}

\begin{proof}
The inequality follows from $\pi_t\geq0$ in
\eqref{eq:ch3-shifted-battery} and $U^*\leq\overline A$ in \eqref{eq:A1}
for every realization. Prices and arrivals are bounded, so the conditional
expectation is well defined. Since $\pi_t$ is determined by $\mathcal F_t$,
we can take it outside the conditional expectation; then
\eqref{eq:energy:conditional-mean} gives the final equality:
\[
 \expect{\pi_t(\overline A-A_{t+1})\mid\mathcal F_t}
 =\pi_t\bigl(\overline A-\expect{A_{t+1}\mid\mathcal F_t}\bigr)=0.\qedhere
\]
\end{proof}

For $\mathbf X^*=U^*\mathbf P^*\in\mathcal A$ as in
\eqref{eq:energy:split-identity}, define the comparison potential
\begin{equation}
\label{eq:ch3-comparison-potential}
  \mathcal H_t(\mathbf{X}^*)
  \coloneqq\frac{(U^*-U_t)^2}{2\eta}
   +\frac{U^*}{\lambda}D_\Psi(\mathbf{P}^*,\mathbf{P}_t)
   +\frac{\pi_t^2}{2\gamma}.
\end{equation}

\begin{Thm:Lemma}[One-step comparison]
\label{lem:energy:combined-potential}
Simultaneously for every $\mathbf X^*\in\mathcal A$, on every sample path,
\begin{equation}
\label{eq:ch3-combined-step}
 c_t(\mathbf X_t)-c_t(\mathbf X^*) +\mathcal H_{t+1}(\mathbf X^*)-\mathcal H_t(\mathbf X^*)
   \leq{}\frac{(\eta+\lambda U^*)G^2+\gamma}{2} +\pi_t(\overline A-A_{t+1}).
\end{equation}
\end{Thm:Lemma}

\begin{proof}
Combining the learning and price bounds gives
\[
 c_t(\mathbf X_t)-c_t(\mathbf X^*) +\mathcal H_{t+1}(\mathbf X^*)-\mathcal H_t(\mathbf X^*)
   \overset{\mathrm{(A)}}{\leq} \frac{(\eta+\lambda U^*)G^2+\gamma}{2} +\pi_t(U^*-A_{t+1}).
\]
For (A), substitute \eqref{eq PArt1}, multiplied by $U^*\geq0$,
and \eqref{eq Part2} into \eqref{eq:energy:split-bound}, and add
\eqref{eq:include}. The potential is defined in
\eqref{eq:ch3-comparison-potential}. The sample-path inequality of
Lemma~\ref{lem:ch3-mean-energy} now gives \eqref{eq:ch3-combined-step}.
\end{proof}

\subsection{Regret bound}

\begin{Thm:Theorem}[Expected regret]
\label{Main Theorem}
Fix positive parameters $\eta$, $\lambda$, and $\gamma$. Suppose the
subgradient and energy bounds in Section~\ref{Assumptions} and the
conditional-mean assumption \eqref{eq:energy:conditional-mean} hold.
Choose $B_{\mathrm{max}}\geq B_{\mathrm{req}}(r)$ from \eqref{eq:Bmax-a}
for some $r\geq0$.
Initialize Algorithm~\ref{Alg1} with $B_1=0$,
$U_1=0$, and $\mathbf{P}_1=(1/n,\ldots,1/n)^\top$.
Then, for every $T\geq1$,
\begin{equation}
\label{eq: Regret}
  \E[\Reg_T(\mathcal A)]
  \leq\frac{(\eta+\lambda \overline A)G^2+\gamma}{2}\,T
       +\frac{\overline A^2}{2\eta}
       +\frac{\overline A\log n}{\lambda}
       +\frac{\gamma B_{\mathrm{max}}^2}{2}.
\end{equation}
\end{Thm:Theorem}

\begin{proof}
Fix a sample path and any comparison allocation
$\mathbf X^*=U^*\mathbf P^*\in\mathcal A$. Uniform initialization and
\eqref{eq:energy:entropy-divergence} give
\[
 D_\Psi(\mathbf P^*,\mathbf P_1)
 =\log n+\sum_{i=1}^n P_i^*\log P_i^*\leq\log n,
\]
because $0\leq P_i^*\leq1$, with $0\log0=0$. Consequently,
\begin{equation}\label{eq:energy:initial-comparison-potential}
\begin{aligned}
 \mathcal H_1(\mathbf X^*)
 &\overset{\mathrm{(A)}}{\leq}\frac{\overline A^2}{2\eta}
   +\frac{\overline A\log n}{\lambda}+\frac{\pi_1^2}{2\gamma}\\
 &\overset{\mathrm{(B)}}{\leq}\frac{\overline A^2}{2\eta}
   +\frac{\overline A\log n}{\lambda}+\frac{\gamma B_{\mathrm{max}}^2}{2}.
\end{aligned}
\end{equation}
In (A), substitute $U_1=0$ in
\eqref{eq:ch3-comparison-potential}, use
$0\leq U^*\leq\overline A$
from \eqref{eq:energy:comparison-set}, and apply the entropy bound above;
$\eta,\lambda>0$ preserve the inequalities. Step (B) uses
$0\leq\pi_1\leq\gamma B_{\mathrm{max}}$ from
\eqref{eq:ch3-shifted-battery} and $\gamma>0$.
Summing \eqref{eq:ch3-combined-step} over slots $1,\ldots,T$ cancels the
intermediate potential terms. The remaining terminal potential is
nonnegative because $U^*\geq0$ and KL divergence is nonnegative
(Lemma~\ref{lem:intro:entropy}). Dropping its negative contribution and
using the initial-potential bound gives
\begin{equation}
\label{main after expet}
\begin{aligned}
 \sum_{t=1}^T\bigl[c_t(\mathbf X_t)-c_t(\mathbf X^*)\bigr]
 &\overset{\mathrm{(A)}}{\leq}
   \frac{(\eta+\lambda\overline A)G^2+\gamma}{2}\,T
   +\frac{\overline A^2}{2\eta}\\
 &\quad+\frac{\overline A\log n}{\lambda}
   +\frac{\gamma B_{\mathrm{max}}^2}{2}
   +\sum_{t=1}^T\pi_t(\overline A-A_{t+1}).
\end{aligned}
\end{equation}
Step (A) uses $U^*\leq\overline A$ in the summed error term and
\eqref{eq:energy:initial-comparison-potential} for the initial potential.
The right side is independent of the comparison allocation. Maximizing
the left side over $\mathbf X^*\in\mathcal A$ therefore gives
$\Reg_T(\mathcal A)$, including when the best allocation depends on the
full realized cost sequence.

The absolute value of each cost difference is at most $2G$
by \eqref{eq:energy:lipschitz}, since both allocations are nonnegative
and have $\ell_1$ norm at most $1$ by \eqref{eq:X-set}.
Thus $|\Reg_T(\mathcal A)|\leq2GT$. The arrival bounds and
\eqref{eq:ch3-shifted-battery} also bound each price term by
$|\pi_t(\overline A-A_{t+1})|\leq\gamma B_{\mathrm{max}}$.
We can therefore take expectations in the resulting regret bound.
By the tower property and \eqref{quu1 bound}, every summand in the price
term has mean zero, giving \eqref{eq: Regret}.
\end{proof}

\subsection{Choosing the parameters}

Fix a horizon $T\geq1$ before running the controller. The following
parameter choice gives the stated regret rate without requiring the mean
arrival $\overline A$. Set the usage step size and price scale to
\begin{equation}
\label{eq:ch3-parameter-scaling}
  \eta\coloneqq\frac{1}{G\sqrt T},\qquad
  \gamma\coloneqq\frac{G}{\sqrt T},\qquad
  \eta\gamma=\frac{1}{T}.
\end{equation}
For $n\geq2$, set the proportion step size to
\begin{equation}
\label{eq:lambda-opt}
  \lambda^*\coloneqq\frac{\sqrt{2\log n}}{G\sqrt T}.
\end{equation}

Choose the sufficient battery capacity
\begin{equation}
\label{eq:ch3-capacity-scaling}
  B_{\mathrm{max}}\coloneqq3\sqrt T.
\end{equation}
Since $T\geq1$, these choices satisfy $0<\eta\gamma\leq1$.
Equation~\eqref{eq:battery-optimal-a} gives $r^*=\sqrt T-1\geq0$,
and the first branch of \eqref{B^*max} gives the capacity above:
$G/\gamma+2/\sqrt{\eta\gamma}=\sqrt T+2\sqrt T=3\sqrt T$.

The choice $\lambda=\lambda^*$ balances the two proportion terms in
\eqref{eq: Regret}: $\lambda^*G^2T/2=\log n/\lambda^*$.
Their sum is therefore
\begin{equation}\label{eq:energy:proportion-tuned-value}
  \overline A\left(\frac{\lambda^*G^2T}{2}+\frac{\log n}{\lambda^*}\right)
  =\overline A G\sqrt{2T\log n}.
\end{equation}
Substituting these choices into \eqref{eq: Regret} gives
\begin{equation}
\label{eq:ch3-explicit-sqrt-bound}
\begin{aligned}
 \E[\Reg_T(\mathcal A)]
 &\overset{\mathrm{(A)}}{\leq}G\sqrt T\left[
   \frac{11+\overline A^2}{2}+\overline A\sqrt{2\log n}\right]\\
 &\overset{\mathrm{(B)}}{\leq}G\sqrt T\left[6+\sqrt{2\log n}\right].
\end{aligned}
\end{equation}
Step (A) uses \eqref{eq:ch3-parameter-scaling} in the usage and price
terms, \eqref{eq:energy:proportion-tuned-value} for the proportion terms,
and $B_{\mathrm{max}}^2=9T$ from \eqref{eq:ch3-capacity-scaling}.
Step (B) uses $0\leq\overline A\leq1$. With $G$ fixed,
\[
  \E[\Reg_T(\mathcal A)]=O\!\left(\sqrt{T(1+\log n)}\right),\qquad
 B_{\mathrm{max}}=O(\sqrt T).
\]

\section{Discussion and conclusion}
\label{sec:energy:conclusion}

The regret guarantee compares the controller with the best fixed
nonnegative allocation that uses no more than the mean energy arrival per
slot. This allocation may require more energy than is available on some
slots. Chapter~\ref{Section:ChapControl} studies
policies that assign fractions of each arrival for use in later slots.
Each comparison policy respects its own energy availability. Its cost
is evaluated using the battery levels and allocations produced by its
decisions.
